% IEEE GRSL letter -- GenCP validation study
% Build: pdflatex letter && bibtex letter && pdflatex letter && pdflatex letter
% Numbers come from ../EVIDENCE.md only. No value enters this file without a row there.

\documentclass[journal]{IEEEtran}

\usepackage{graphicx}
\usepackage{amsmath}
\usepackage{booktabs}
\usepackage[hidelinks]{hyperref}

\graphicspath{{../figures/}}

\begin{document}

% Title changed 2026-08-26. The previous title, "Adversarial training degrades the
% geometric accuracy of synthetic reference imagery", was narrower than the result:
% C5 carries no discriminator anywhere in its objective and hallucinates hardest of
% the five arms, so the mechanism is plausibility pressure, not the adversarial term.
% This is the working title from letter-skeleton.md; letter.tex had never caught up.
\title{Plausibility Pressure Degrades Generated Reference Imagery
for Geometric Matching}

\author{Mustafa Vedat Y\i ld\i r\i m and co-author%
\thanks{Manuscript submitted 2026. Study repository and full experiment records:
\url{https://github.com/mvy0502/gencp-validation}.}}

\markboth{IEEE Geoscience and Remote Sensing Letters}%
{Y\i ld\i r\i m: Plausibility pressure and generated reference geometry}

\maketitle

\begin{abstract}
% ABSTRACT -- 200 words, and it is the tightest-governed text in the letter.
% Source: letter-skeleton.md, "Title block and abstract"; revised 2026-08-26.
%
% Required content, IN THIS ORDER:
%   1. Setting: generated reference imagery for satellite georeferencing.
%   2. Design: 2x2 factorial crossing an adversarial term with L1 vs LPIPS
%      reconstruction, everything else held fixed.
%   3. Primary number -- the SEED-LEVEL six-seed sign replication:
%      adversarial-off beats adversarial-on under LPIPS in all six confirmatory
%      seeds, P = 1/64, direction fixed in advance.
%   4. Generalisation: the effect replicates under both reconstruction terms, and
%      LPIPS alone invents more than either adversarial arm.
%   5. Mechanism: edge density in input-silent regions.
%   6. The design rule.
%
% TWO HARD RULES.
% (i) The abstract carries NO INTERVAL AT ALL. Not the six-seed interval, not any
%     other. A sign replication with its P-value is the claim. An interval in an
%     abstract invites the chip-level-as-treatment-level reading that this whole
%     seed-replication package exists to correct. Results III.2 may report the
%     interval; the abstract may not.
% (ii) The old chip-level primary -- C5-C4 = -0.487 +- 0.053 px, t = -9.2, better on
%      113/130 chips -- is STRUCK and must not appear here. It measures how
%      consistently one checkpoint beats another across 130 chips, not how
%      consistently the treatment works. That conflation is the paper's own
%      methodological point.
%
% Index terms: image registration, generative adversarial networks, ground control
% points, perceptual loss, georeferencing, Sentinel-2.

% P4 block 3, 2026-09-13: four claims brought down to the body (LR exception; comparison
% class of "worse"; "necessary condition" not "explains"); the L1-family reading added to
% III-B so "and under L1 as well" has an explicit body statement. Rule: the abstract never
% speaks more strongly than the body.
% DRAFTED 2026-09-13. Order per the comment block above; NO interval; the struck chip-level
% primary does not appear. Word count checked by script at commit time.

Generated reference imagery has been proposed for satellite georeferencing: a conditional
GAN renders a Sentinel-2-like image from OpenStreetMap and land cover. A
$2\times2$ factorial crosses an adversarial term with the reconstruction loss, L1 against the perceptual loss LPIPS, everything else held fixed except a disclosed learning-rate warm-up on the adversarial arms whose consequence is bounded, and scores
each arm by the
positional error of feature matches against real Sentinel-2 over 130 chips. Removing the
adversarial term improves positional accuracy under the perceptual loss in all six
confirmatory training seeds ($P = 1/64$, direction fixed in advance), and under L1 as well;
and the perceptual loss alone, with no discriminator anywhere, carries its own positional
penalty in all six seeds. An input-conditioned measurement supplies the necessary condition, not a complete
explanation: the edge density
an arm renders where the conditioning input asserts no structure is lowest for the L1-only
arm and highest for the LPIPS-only arm, which invents more than either adversarial arm. The mechanism is plausibility pressure, not the adversarial term. Matched directly, where it matches, the rasterised map itself localises better than every
generated arm, including at equal counts; what remains of the case for a generated
reference is point yield and coverage. The design rule: if you generate a reference image for a geometric consumer, do not
train it under plausibility pressure.

\end{abstract}

\begin{IEEEkeywords}
Geometric calibration, ground control points, image-to-image translation,
generative adversarial networks, Sentinel-2, OpenStreetMap.
\end{IEEEkeywords}

\section{Introduction}
% I. INTRODUCTION -- 600 words. NOW ABSORBS THE FORMER 02-scope.tex.
%
% What a GCP reference is for, why synthetic references were proposed, and the
% question this letter answers: if you generate one, how should you train it?
%
% SCOPE, folded in here as TWO SENTENCES, not a section and not a result.
% (paper-roadmap.md:152 -- "E1/E2 become two sentences of scope in the
% introduction, not results." The former 02-scope.tex is deleted, its content lives
% here.)
%   E1: no availability gap -- 0/24 extents without a usable scene, median freshness
%       2 days, max 17.
%   E2: no currency advantage -- ABSENT, +0.008 +- 0.031 px.
%   Footnote: scene-per-year counts are right-censored lower bounds (archive query
%       capped at 100, every extent hit the cap); the censoring touches neither gap
%       rate nor freshness. State honestly that E2's null is equally consistent with
%       "currency does not help" and "OSM had not yet recorded the change"
%       (open-items.md item 21).
%
% E3 DOES NOT APPEAR -- not here, not anywhere in the letter (corrections-log
% entry 16, reclassified exploratory). It moved to the second paper.
%
% Related work: compressed to ~150 words and 4 citations, each with an explicit
% non-negotiable justification. Do not let this grow back to 7.

TODO.


\section{Materials and Methods}
% II. MATERIALS AND METHODS -- 750 words. (Was 03-method.tex.)
%
% Contents:
%   - The 2x2 loss factorial: GAN x L1/LPIPS, everything else held fixed.
%   - The evaluation instrument: KARIOS. Config karios_gencp.json is sha256-pinned
%     in the registrations.
%   - Reference construction.
%   - Registration discipline: predictions and falsification bands committed before
%     each run. This is the letter's strongest procedural claim -- the factorial was
%     registered at b07e719 before any run, harness committed at 40cde9b, Gates 0/1
%     passed, every registered band fired, retraction condition not triggered.
%   - The 1/256 scale bug: ONE PARAGRAPH plus a repo pointer. It is a methods note
%     here, not a result. (+0.390625% = exactly 1/256; true GSD 10.0390625 m vs 10.0
%     declared; worst case 14.1 m at the SE corner, zero at NW.)
%   - Seed defence, one sentence: pre-registered n >= 60 single-draw rule, regA
%     det/stoch bound <= 0.05 px, measured effect sizes 0.38-0.70 px.
%
% METHODS/RESULTS SPLIT (Resolution 2): empirical results do NOT live in this
% section. If a number is being reported rather than a procedure described, it
% belongs in III.

TODO.


\section{Results}
% III. RESULTS -- 800 words, Table I, Fig. 1. (Was 04-results.tex.)
%
% THE GUIDANCE THAT WAS HERE IS SUPERSEDED. The old comments named T1's C1 row as
% the lead measurement and cited "mediation 0%". Both are dead:
%   - paper-roadmap.md:180 -- "the factorial replaces T1's C1 row as the primary
%     measurement". T1 moved to the second paper along with the ODTU/Cappadocia
%     contamination pair and E3.
%   - B3 part 2 (mediation) "does not appear at all -- it is void as stated"
%     (B2-B3-audit.md, corrections-log entry 20). The reported conditional was the
%     OLS fitted value at the covariate means, algebraically the raw mean, so
%     "loses 0% of its magnitude" could not have been a mediation result.
%
% Structure, per letter-skeleton.md III (subsection: words):
%   1. The panel, 80. Table I. Ordering C2 < C5 < C4 < C1 < pretrained.
%   2. Primary, 120. SEED-LEVEL: C5-C4 negative in all six confirmatory seeds
%      (45-50), P = 1/64, direction fixed in advance by seed 42. Per-seed:
%      -0.6153, -0.6462, -0.6162, -0.5942, -0.6054, -0.5775 px. Seed-level mean
%      -0.609 px, 95% interval [-0.634, -0.585] (df = 5, t* = 2.571) -- REPORTED,
%      NOT REQUIRED; the registered reading is the sign replication and only that.
%      Seed 42's -0.487 is the generating observation and falls OUTSIDE the six-seed
%      range: do not present it as one of the replicates.
%   3. Interaction -- DEAD, REPLACED IN FULL, NOT AMENDED. The registered seed-level
%      interaction failed at 5/6 on all three pre-specified scales (raw, log,
%      within-chip rank), the same seed (46) breaking each. Pre-committed
%      consequence: no interaction claim; "substitutes" and "the same lever" are
%      removed from the paper. The 120 words do NOT return to the budget -- they buy
%      the mandatory disclosure (paper-context-addendum.md §24). Five elements may
%      not be lost when shortening: registered in advance, all three scales, 5/6
%      with the same seed breaking each, no claim made, the other block reported
%      with its weight stated. NEVER write "the interval excludes zero, so the
%      interaction is real" or anything functioning as it.
%   4. Secondary, 60. C5-C2 positive in all six seeds, P = 1/64: perceptual
%      reconstruction carries its own positional penalty with NO discriminator
%      present. Mean +0.063 px, interval [+0.027, +0.098], reported not required.
%      STATE the thinnest margin, +0.0068 px at seed 46. This subsection is what
%      kept the title: had it failed in any seed, the claim would have narrowed from
%      "plausibility pressure" back to "the adversarial term".
%   5. Dose-response, 80 + Fig. 2. Penalty by epoch under LPIPS: 0.334 -> 0.254 ->
%      0.441 -> 0.496 -> 0.487, all >= 6 SE. Same dip-then-grow-then-plateau shape
%      as the L1 family. B1 is SUPPORT ONLY and must not be presented as the spine:
%      it is post-hoc, the registered bands did not cover the observed shape.
%   6. Mechanism, 140 + Fig. 1. The edge-ratio column: C2 0.28 vs C1 1.10 / C4 1.12
%      / C5 1.16, pretrained ~1.02. With no discriminator anywhere, LPIPS alone
%      invents most -- the registered primary prediction for C5 and the single
%      result that widens the claim from "adversarial" to "plausibility pressure".
%      CARRY the section-22 non-monotonicity caveat (paper-context-addendum.md §22):
%      the ratio separates the restrained arm from the unrestrained ones but does
%      NOT order the errors within the unrestrained group. Invention is a necessary
%      condition, not a complete explanation.
%   7. Point-count argument, 160 (100 + 60). The objection: L1-only just produces
%      fewer features, and fewer-but-better is a trivial trade-off. C5 refutes it ON
%      THIS PANEL, and the panel is named in the sentence: C5 produces more
%      surviving matches than C2 (median 88 vs 72) and still scores worse. DO NOT
%      cite the production-path point counts here -- they reverse.
%
% B2 stays as the production-path ablation row.

TODO.


\section{Alternative Explanations}
% IV. ALTERNATIVE EXPLANATIONS -- 500 words, Table II. NEW FILE.
%
% A full-length draft already exists in this repository at
% manuscript/drafts/draft-section-IV.md (moved from the study repository's
% tubitak/docs/ on 2026-09-13; 1,109 words, four refutations plus the
% design-based argument). It is drafted at arXiv length by the structural decision
% of 2026-08-26; condensing it to the letter's 500 words is a separate task and has
% not been done. DO NOT paste the long version in here.
%
% TABLE II IS THREE ROWS, not five, and not the four of the first revision:
%   1. Cold-discriminator.
%   2. Matcher independence -- re-attributed across two registrations. The "48 of
%      49" figure is STRUCK and must not be quoted: it is not reproducible from the
%      artifact.
%   3. Blur vs restraint -- REPLACED by a stronger test. The original was a
%      single-seed negative control; the replacement is the six-seed informative-
%      mask edge-ratio measurement. C2 reproduces the real image's edge density to
%      within 1.5% where the input asserts structure, so the suppression is
%      SELECTIVE (restraint), not blur. Verdict confirmed across all six seeds.
%
% REMOVED ROWS, and they do not come back:
%   - Mediation row: STRUCK (corrections-log entry 20; paper-roadmap already rules
%     B3 part 2 does not appear at all).
%   - Georeferencing row: REMOVED from the table and ANSWERED IN PROSE instead. The
%     answer is design-based rather than measured -- stronger than the 86% figure it
%     replaces, because it depends on the design rather than on one number. The
%     registered positional contrasts contain no pretrained comparisons, so the
%     objection does not reach them.
%
% The freed word budget is NOT reclaimed by this section; the allocation stands as
% revised.
%
% Mediation footnote: write it carefully. The registered mediation test was run and
% its reported conditional is void as stated -- say that, do not imply the test
% supports anything.
%
% Mechanistic note, ~40 words, because it explains why the result is
% matcher-independent.

% Migrated from manuscript/drafts/draft-section-IV.md (Draft 1, 2026-08-26, arXiv length,
% 1,109 words) on 2026-09-13. Prose only, transcribed without compression or correction;
% the condensation to the letter's 500 words is a separate task and has NOT been done.
% Markdown bold is carried as \textbf. Table II is not typeset; the reference is left as a
% placeholder. Values that are the Table I edge-ratio column on the input-silent mask are
% wrapped in \tabone{...}.

Four explanations compete with the account above. Three are tested and refuted; the
fourth is answered by the design itself and needs no test.
Table~\todo{II: three rows, not yet typeset} summarises; the prose gives each row the
evidence that supports it and the caveat that qualifies it.

\subsection{It is blur, not restraint}

\textbf{The objection.} The L1-only arm produces smoother output, and feature matchers
prefer smoother imagery. If so, its advantage is a mechanical consequence of blur rather
than of any learned decision about what to render.

\textbf{The test.} The invention measurement is defined on input-silent pixels---those
where the conditioning input asserts no structure. Computing the \emph{same} ratio on the
complementary mask, where the input does assert structure, separates the two explanations
exactly, because \textbf{blur suppresses edges uniformly while restraint suppresses them
conditionally.} A smoothing explanation predicts the L1-only arm below unity on both
masks. A restraint explanation predicts it near unity where the input is informative and
far below where it is not. The thresholds separating ``near unity'' from ``suppressed''
were fixed in advance at 0.80 and 0.50, reusing the bands already registered for the
original measurement rather than choosing new ones.

\textbf{The result.} Across six seeds the L1-only arm reproduces the real image's edge
density to within 1.5\% on the informative mask---0.986, range 0.980 to 0.989---while
sitting at \tabone{0.277} on the input-silent mask. \textbf{A factor of 3.6 between the two
masks, in the same arm, on the same chips, with the same operator, in every seed.} Both
conditions hold six times out of six. The other three arms sit between 1.03 and 1.05 on
the informative mask, slightly above reality, consistent with the \tabone{1.07 to 1.16}
they show where the input says nothing.

\textbf{The reading.} Every arm except the L1-only one adds edges everywhere. The L1-only
arm adds them only where the input warrants it. \textbf{That is conditional suppression,
which blur cannot produce}, and the objection is refuted by a positive test rather than by
a negative control.

\textbf{What this row is not.} It does not establish that the L1-only arm is optimal, nor
that restraint is the only mechanism in play---Section III's honest-limit paragraph
stands. It establishes that the smoothing explanation, specifically, does not fit the
data.

\subsection{Cold-started discriminator damage}

\textbf{The objection.} The published generator is deposited without its discriminator,
so every adversarial arm here begins from a randomly initialised one. Adversarial arms
might therefore lose because of a damaged start rather than because of the adversarial
objective.

\textbf{The test.} A checkpoint sweep at epochs 1, 2, 5, 10 and 20, comparing each arm to
the pretrained generator and to its own counterpart.

\textbf{The result.} At epoch 1 the adversarial+L1 arm is \textbf{already better than
pretrained}, by $-0.399 \pm 0.064$~px at 6.3 standard errors. \textbf{That is the wrong
sign for startup damage}: an arm crippled by a cold discriminator would begin worse than
the checkpoint it started from, not better. The adversarial deficit relative to the
non-adversarial arm is present from epoch 1 at $+0.546 \pm 0.048$ and reaches $+0.700$ by
epoch 20.

\textbf{The caveat, and it is ours rather than a reviewer's.} The path from epoch 1 to
epoch 20 is not monotone: it dips to $+0.384$ at epoch 5 before growing. \textbf{The
registered bands for this sweep did not anticipate that shape, and the conclusion is read
off the curve rather than matched to a band.} We state it that way because presenting it
as a clean band hit would misrepresent how it was obtained. The refutation does not depend
on the shape---it depends on the epoch-1 magnitude, which is large and correctly signed.

\subsection{Optimising the evaluation metric}

\textbf{The objection.} The result might be an artefact of the KLT matcher used to
produce it.

\textbf{The evidence, from two independent registrations.} The first varied the
\emph{descriptor family}: ORB, AKAZE and mutual information all rank L1-only ahead of
adversarial+L1 on both the Ankara and European sets, with ORB at $-0.613 \pm 0.135$~px.
\textbf{Two caveats travel with these numbers and are stated rather than buried.} The ORB
figure is a paired difference over the \textbf{29 chips where both arms matched}, not over
the 53 chips where the L1-only arm matched at all; the AKAZE figure rests on 11 paired
chips. And the mutual-information margin of $-1.260 \pm 0.261$ is a \textbf{lower bound},
not an estimate: the registered subpixel refinement never ran, and the search grid censors
15.8\% of chips at its bound, censoring the worse arms harder and therefore compressing the
margin toward zero.

The second registration varied the \emph{matcher family} rather than the descriptor: KLT,
an NCC template grid, and phase correlation, crossed with two band conversions and an
urban subset over 6,510 scored comparisons. \textbf{The arm ordering is preserved in every
condition cell except two}, both at the European urban subset under phase correlation,
one in each band conversion, and both far below the threshold at which an ordering change
would be claimed. A registered prediction that the template matcher would favour the
sharper arm \textbf{failed in the direction that strengthens the result}: the L1-only
arm's margin grew rather than shrank, in all eight set-by-conversion combinations.

\textbf{The reading.} The candidate is refuted by two registrations that share no matcher,
no chip set and no scoring code. \textbf{That is stronger evidence than either alone}, and
it is the reason this row is reported as two tests rather than one.

\subsection{Corrected georeferencing in the fine-tuning pairs}

\textbf{The objection.} The fine-tuning pairs may be better georeferenced than the data
the published generator saw, so the advantage would come from the training data rather
than from the loss.

\textbf{The answer is in the design and needs no measurement.} No registered positional
contrast compares a fine-tuned arm with the pretrained generator: all four arms are
fine-tuned on identical pairs, so any georeferencing improvement is common to them and
cancels.

\textbf{Why this replaces a measurement we previously reported.} An earlier decomposition
attributed roughly 86\% of the European gain to scatter reduction rather than to a
systematic shift, and was reported as refuting this candidate. \textbf{That measurement's
per-chip artifacts did not survive and it cannot be reproduced}, which is disclosed in the
data-availability statement. The design argument is not a fallback: \textbf{it is
stronger, because it depends on the structure of the experiment rather than on a
computation, and nothing can be lost that would make it unverifiable.}

\subsection{Why the result is matcher-independent, mechanically}

A blurred template gives a broad correlation peak; invented structure gives a sharp peak
in the wrong place. \textbf{A broad peak in the right place localises better than a sharp
peak in the wrong one}, and that is true of any matcher that localises by correlation.
This is offered as the mechanical reason the previous row's result is unsurprising, not
as further evidence for it.


\section{Discussion, Limitations, and the Design Rule}
% V. DISCUSSION, LIMITATIONS, DESIGN RULE -- 450 words.
% (Was 06-discussion.tex. The former 07-conclusion.tex is MERGED IN HERE -- there is
% no separate conclusion section in the current five-section structure.)
%
% Structure, per letter-skeleton.md V (subsection: words):
%   1. The design rule, 60. If you generate a reference image for a geometric
%      consumer, do not train it under plausibility pressure. This is the sentence
%      the whole letter exists to license -- it also serves the conclusion role the
%      deleted 07-conclusion.tex used to hold.
%   2. Relation to the theory, 90. An empirical instance of the
%      perception-distortion trade-off.
%   3. The Cramer-Rao pre-emption, ONE sentence. The classical result says variance
%      scales -- pre-empt the reviewer who raises it.
%   4. Limitations, 180, compressed. The institution's own matching software was
%      never run against this. Carry the honest caveats that must travel with each
%      leg, including the E2 ambiguity (open-items 21) and the section-22
%      non-monotonicity caveat on the mechanism.
%   5. What generalises, 80. NOT "GANs are bad". The claim is about the CONSUMER:
%      any geometric consumer of generated imagery. Scope boundary in one sentence
%      -- the 10 m premise fails; sub-metre is where it binds.
%   Plus, from the old conclusion: the pointer to the open study repository.
%
% E3 does not enter a results table. If it appears at all it is labelled exploratory
% and only in the arXiv long version -- not here.

TODO.


\bibliographystyle{IEEEtran}
\bibliography{refs}

\end{document}
